%! Choosing and Comparing Models in Context — Unit 2, Lesson 2.7 — Guided Notes (Answer Key)
\documentclass[10pt]{article}
\usepackage{atda-article}
\usepackage{atda-key}   % pulls in -boxes; do NOT also load -boxes
\newcommand{\TallMath}[1]{$\displaystyle #1\rule[-1.4em]{0pt}{3.2em}$}
\setlength{\workrowsep}{6pt}

\begin{document}

\pageheader{Unit 2, Lesson 2.7}{Guided Notes --- Answer Key}
% NAMESTRIP: no \namedateperiod here — mirrors the blank. Teacher-only prose goes
% in the lesson plan as \begin{teachernote}[Guided Notes], never in a key.

\vspace{0.15in}

\begin{objectivebox}
\small
By the end of this lesson, I will be able to\dots
\begin{enumerate}[leftmargin=1.5em, itemsep=2pt, topsep=3pt]
  \item decide whether a \textbf{linear}, \textbf{quadratic}, or \textbf{exponential} function best
        models a table, a graph, or a story --- and say which test decided it;
  \item run the two table tests in order --- \textbf{first differences}, then \textbf{second
        differences}, then \textbf{ratios} --- and know that they only work on \emph{equal} steps in
        the input;
  \item compare and contrast the three families across every representation I have: parent function,
        table, graph, characteristics, and the story itself;
  \item use a model to answer a question in context, \textbf{graphically} and \textbf{algebraically};
  \item say what a model \emph{cannot} tell me --- where its answers stop being trustworthy, and why.
\end{enumerate}
\end{objectivebox}

\vspace{0.10in}

\boxguard
\begin{vocabbox}
\small
Fill in each term as we build it together. The first one is the word this whole lesson is about.
\vocabans{Mathematical model}{a function chosen to describe a real situation. It is a
    \emph{choice}, not a fact --- a different choice can fit the same data, and a model can be
    wrong.}
\vocabans{First differences}{the changes from one output to the next, when the inputs go up in equal
    steps. Constant first differences mean the model is \textbf{linear}.}
\vocabans{Second differences}{the changes in the first differences. Constant second differences
    (when the first ones are not constant) mean the model is \textbf{quadratic}.}
\vocabans{Common ratio}{the number each output is multiplied by to get the next one, when the inputs
    go up in equal steps. A constant ratio means the model is \textbf{exponential}.}
\vocabans{Extrapolation}{using a model to predict beyond the inputs the data actually covered. It is
    where a wrong choice of family shows up --- and where even a right one stops being believable.}
\end{vocabbox}

\vspace{0.10in}

\boxguard[12]
\begin{hookbox}
\small
\textbf{Two students, one table, and both of them are right about what they checked.} A school posts
a video on Monday. Here are the views, in thousands, on each of the first four days.

\smallskip
\begin{center}
\renewcommand{\arraystretch}{1.25}
\begin{tabular}{@{}|l|c|c|c|c|@{}}
\hline
Day $d$ & $0$ & $1$ & $2$ & $3$ \\ \hline
Views $V$ (thousands) & $3$ & $6$ & $12$ & $24$ \\
\hline
\end{tabular}
\end{center}

\smallskip
Student A: \textbf{``The views go up by $3$, then by $6$, then by $12$. The jumps keep getting
bigger, and a parabola's jumps get bigger too --- so this is \emph{quadratic}.''}

\smallskip
Student B: \textbf{``Every day is double the day before --- so this is \emph{exponential}.''}

\smallskip
\textbf{Everything both students said about these numbers is true.} The jumps really do get bigger,
and each value really is double the one before. Decide which family models the video, and say what
Student A checked that was not enough.

\ansline{Student B. The jumps $3, 6, 12$ are not growing by a constant amount --- they are \emph{doubling} --- so the}\\
\ansline{second differences are not constant and the model is not quadratic. Growing jumps only rule out linear.}\\
\end{hookbox}

\vspace{0.10in}

\boxguard[30]
\begin{notesbox}{1. Two Tests, Run in Order}
\small
Choosing a family is a \textbf{decision}, and you can decide wrong. When you have a table, two tests
make the decision for you --- and you run them in a fixed order.

\vspace{4pt}
For each table below, fill the three rows: the \textbf{change} from one output to the next, the
\textbf{change in those changes}, and the \textbf{ratio} of each output to the one before it.

\vspace{6pt}
\textbf{(a) A rideshare fare.} $F$ is the fare in dollars for a ride of $m$ miles.

\vspace{2pt}
\renewcommand{\arraystretch}{1.3}
\begin{tabular}{@{}|l|c|c|c|c|@{}}
\hline
$m$ (miles) & $0$ & $1$ & $2$ & $3$ \\ \hline
$F$ (\$) & $3$ & $5$ & $7$ & $9$ \\ \hline
First differences & --- & \ans{$+2$} & \ans{$+2$} & \ans{$+2$} \\ \hline
Second differences & --- & --- & \ans{$0$} & \ans{$0$} \\ \hline
Ratios & --- & \ans{$\approx 1.67$} & \ans{$1.4$} & \ans{$\approx 1.29$} \\
\hline
\end{tabular}

\vspace{6pt}
\textbf{(b) Braking distance.} $D$ is the distance in feet a car travels after the brakes are hit,
at a speed of $v$ miles per hour. \emph{The input steps are $10$, not $1$ --- that is fine, as long
as they are all the same.}

\vspace{2pt}
\renewcommand{\arraystretch}{1.3}
\begin{tabular}{@{}|l|c|c|c|c|@{}}
\hline
$v$ (mi/h) & $20$ & $30$ & $40$ & $50$ \\ \hline
$D$ (feet) & $20$ & $45$ & $80$ & $125$ \\ \hline
First differences & --- & \ans{$+25$} & \ans{$+35$} & \ans{$+45$} \\ \hline
Second differences & --- & --- & \ans{$+10$} & \ans{$+10$} \\ \hline
Ratios & --- & \ans{$2.25$} & \ans{$\approx 1.78$} & \ans{$\approx 1.56$} \\
\hline
\end{tabular}

\vspace{6pt}
\textbf{(c) A shared post.} $P$ is the number of people who have seen a post, $n$ hours after it
went up.

\vspace{2pt}
\renewcommand{\arraystretch}{1.3}
\begin{tabular}{@{}|l|c|c|c|c|@{}}
\hline
$n$ (hours) & $0$ & $1$ & $2$ & $3$ \\ \hline
$P$ (people) & $5$ & $15$ & $45$ & $135$ \\ \hline
First differences & --- & \ans{$+10$} & \ans{$+30$} & \ans{$+90$} \\ \hline
Second differences & --- & --- & \ans{$+20$} & \ans{$+60$} \\ \hline
Ratios & --- & \ans{$3$} & \ans{$3$} & \ans{$3$} \\
\hline
\end{tabular}

\vspace{6pt}
\textbf{Now write the rule.} Exactly one row came out constant in each table.

\vspace{2pt}
\renewcommand{\arraystretch}{1.4}
\begin{tabularx}{\linewidth}{@{}p{6.0cm} X@{}}
\rowcolor{mist}
\textbf{If this row is constant\dots} & \textbf{\dots the model is} \\
\hline
the first differences & \ans{linear --- like the rideshare fare} \\
the second differences (but not the first) & \ans{quadratic --- like the braking distance} \\
the ratios & \ans{exponential --- like the shared post} \\
\end{tabularx}

\vspace{6pt}
\textbf{The one condition, and it is not optional.} Both tests only work when the inputs go up in
\ans{equal} steps. If the table jumps $1, 2, 5, 10$, the differences mean nothing.

\vspace{6pt}
\textbf{Now settle the hook.} Run both tests on the video table.

\vspace{2pt}
\renewcommand{\arraystretch}{1.3}
\begin{tabular}{@{}|l|c|c|c|c|@{}}
\hline
Day $d$ & $0$ & $1$ & $2$ & $3$ \\ \hline
$V$ (thousands) & $3$ & $6$ & $12$ & $24$ \\ \hline
First differences & --- & \ans{$+3$} & \ans{$+6$} & \ans{$+12$} \\ \hline
Second differences & --- & --- & \ans{$+3$} & \ans{$+6$} \\ \hline
Ratios & --- & \ans{$2$} & \ans{$2$} & \ans{$2$} \\
\hline
\end{tabular}

\vspace{6pt}
The second differences are \ans{not constant}, so the model is \emph{not} quadratic. The ratios are all
\ans{$2$}, so the video is modelled by an \ans{exponential} function. \textbf{Student
\ans{B} is right.}

\vspace{4pt}
\textbf{And here is what Student A actually checked.} ``The jumps get bigger'' rules out
\ans{linear} --- and nothing else. \emph{Every} family except that one has growing jumps, so that
observation cannot tell quadratic from exponential. \textbf{You have to check which row is
constant, not just which row is growing.}
\end{notesbox}

\vspace{0.10in}

\boxguard[30]
\begin{notesbox}{2. The Same Three Families, Five Ways}
\small
This is the summary sheet for the whole unit. Fill the empty cells --- everything in them you have
already met.

\vspace{4pt}
\renewcommand{\arraystretch}{1.45}
\begin{tabularx}{\linewidth}{@{}p{2.9cm}|X|X|X@{}}
\rowcolor{mist}
& \textbf{Linear} & \textbf{Quadratic} & \textbf{Exponential} \\
\hline
Parent (2.2) & $f(x)=x$ & $g(x)=x^2$ & $h(x)=b^{\,x}$ \\
\hline
Table test & \ans{first differences} constant & \ans{second differences} constant & \ans{ratios} constant \\
\hline
Graph shape & \ans{a straight line --- it never turns and never bends} & one turning point, and the two sides are mirror images &
    \ans{no turning point; it bends harder and harder at one end} \\
\hline
Increasing / decreasing (2.4) & one or the other, on the whole domain & one of each, meeting at the
    turning point & one or the other, on the whole domain \\
\hline
Absolute max / min (2.4) & neither, unless the domain stops & exactly one, at the
    \ans{turning point} & neither, unless the domain stops \\
\hline
Asymptote (2.5) & none & none & \ans{exactly one, horizontal} \\
\hline
End behavior (2.5) & one end runs up, the other runs down & both ends run the \emph{same} way &
    one end runs off; the other flattens toward the asymptote \\
\hline
What the story says & the same \ans{amount} is added each step & something rises and then falls
    (or falls and then rises) --- there is a turnaround & the same \ans{factor} multiplies each
    step \\
\end{tabularx}

\vspace{6pt}
\textbf{The two rows that do the most work.} A story that says ``\emph{plus \$30 a month}'' and a
story that says ``\emph{up $30\%$ a month}'' sound almost the same out loud and are two different
families --- the first \ans{linear}, the second \ans{exponential}. And the fastest thing to look for on
a graph is a \textbf{turning point}: if the graph turns around, the model is \ans{quadratic}, and if it
never turns you still have two families left to separate.
\end{notesbox}

\vspace{0.10in}

\boxguard[30]
\begin{notesbox}{3. Using the Model --- and What It Cannot Tell You}
\small
The video is exponential, and the model that fits the table is
\[
  V(d) = 3 \cdot 2^{\,d}, \qquad V \text{ in thousands of views, } d \text{ in days.}
\]

\begin{center}
\begin{tikzpicture}
\begin{axis}[
    width=10.4cm, height=4.6cm,
    xlabel={$d$ (days since the video was posted)}, ylabel={$V$ (thousands of views)},
    xmin=0, xmax=6, ymin=0, ymax=200,
    xtick={0,1,2,3,4,5,6}, ytick={0,25,50,75,100,125,150,175,200},
    minor x tick num=1,
    grid=both, grid style={linegray!45}, minor grid style={linegray!30},
    axis lines=left,
    tick label style={font=\tiny}, label style={font=\scriptsize}]
  \addplot[royal, very thick, domain=0:6, samples=80]{3*2^x};
  \addplot[keyred, only marks, mark=*, mark size=1.7pt] coordinates {(4,48) (5,96)};
  \addplot[keyred, thick, dashed] coordinates {(0,50) (6,50)};
\end{axis}
\end{tikzpicture}
\end{center}

\vspace{-2pt}
\textbf{Algebraically --- how many views on day $5$?} Put $5$ in for $d$:

\vspace{2pt}
\begin{work}
  V(5) &= 3 \cdot 2^{5} \\
       &= 3 \cdot 32 \\
       &= 96 \text{ thousand views}
\end{work}

\vspace{2pt}
\textbf{Graphically --- when does the video first pass $50$ thousand views?} Draw across from $50$ on
the vertical axis until you meet the curve, then read straight down. It happens between day
\ans{$4$} and day \ans{$5$}, because $V(4)=$ \ans{$48$} and $V(5)=$ \ans{$96$}.

\vspace{6pt}
% Unbroken, notes box 3 spilled just its last two lines onto p4 -- a stub with the
% whole box on the page before. \boxguard is inert inside a breakable box, so
% \tcbbreak is the lever; breaking here gives p4 a substantial chunk instead.
% Mirrored in notes_key.
\tcbbreak
\textbf{Now the question that matters most.} Use the model to predict the views on day $20$:
$3 \cdot 2^{20} = 3{,}145{,}728$ thousand, which is about \textbf{$3.1$ billion views}. That is more
views than any video has ever had, from a school video, in three weeks.

\vspace{4pt}
\textbf{The model is not broken --- it is being used outside the days the data covers.} Reading a
model beyond the inputs your data covered is called \ans{extrapolation}, and it is where wrong family
choices show up. Doubling has to stop: the school runs out of \ans{people} to show the video to.

\vspace{4pt}
\textbf{Write down the sentence.} A model is a \ans{choice}, not a fact, and every model has a
\ans{window} of inputs where it works. Say what the model predicts \emph{and} whether you believe it.
\end{notesbox}

\vspace{0.10in}

\boxguard[30]
\begin{practicebox}
\small
Three quantities were measured. Each one is modelled well by one of our three families.

\vspace{4pt}
\renewcommand{\arraystretch}{1.3}
\begin{center}
\begin{tabular}{@{}|l|c|c|c|c|c|@{}}
\hline
\textbf{(a)} $t$ (hours draining) & $0$ & $1$ & $2$ & $3$ & $4$ \\ \hline
Water left (gallons) & $4000$ & $3650$ & $3300$ & $2950$ & $2600$ \\
\hline
\end{tabular}

\smallskip
\begin{tabular}{@{}|l|c|c|c|c|c|@{}}
\hline
\textbf{(b)} $t$ (seconds after the punt) & $0$ & $1$ & $2$ & $3$ & $4$ \\ \hline
Height of the ball (feet) & $0$ & $48$ & $64$ & $48$ & $0$ \\
\hline
\end{tabular}

\smallskip
\begin{tabular}{@{}|l|c|c|c|c|c|@{}}
\hline
\textbf{(c)} $h$ (hours since the dose) & $0$ & $5$ & $10$ & $15$ & $20$ \\ \hline
Caffeine in the blood (mg) & $160$ & $80$ & $40$ & $20$ & $10$ \\
\hline
\end{tabular}
\end{center}

\begin{enumerate}[leftmargin=1.6em, itemsep=6pt, topsep=6pt]
  \item Run the tests on each measurement and fill the table. Name the family, and name the row that
        came out constant.
\smallskip

\renewcommand{\arraystretch}{1.4}
\begin{tabularx}{\linewidth}{@{}p{1.0cm}|X|X|X@{}}
\rowcolor{mist}
& \textbf{Which row is constant?} & \textbf{Its constant value} & \textbf{The family} \\
\hline
(a) & \ans{first differences} & \ans{$-350$ gal/hour} & \ans{linear} \\
(b) & \ans{second differences} & \ans{$-32$ ft} & \ans{quadratic} \\
(c) & \ans{ratios} & \ans{$0.5$ every $5$ h} & \ans{exponential} \\
\end{tabularx}

  \item Use (b). Give the maximum height of the punt and the time it happens, then give the two
        zeros and say what each one means about the football.

\ansline{Maximum height $64$ feet, at $t=2$ seconds --- the turning point, with the two sides mirroring each other.}\\
\ansline{Zeros at $t=0$ and $t=4$ seconds. The zero at $t=0$ is the instant the ball is punted, and the zero}\\
\ansline{at $t=4$ is the instant it hits the ground --- the two times the ball's height is $0$.}\\

  \item Use (c). The dose halves every $5$ hours, so the model is $C(h) = 160 \cdot (0.5)^{h/5}$
        milligrams. How much caffeine is left after $25$ hours? Does this model ever say the caffeine
        reaches \emph{exactly} zero? Answer with a Lesson 2.5 word.

\ansline{$25$ hours is five half-lives, so $C(25) = 160 \cdot (0.5)^5 = 160 \div 32 = 5$ milligrams.}\\
\ansline{No. Half of a positive number is still positive, so the amount never reaches $0$ --- it only gets}\\
\ansline{closer and closer. The line $C=0$ is a \textbf{horizontal asymptote}: approaching is not arriving.}\\

  \item Both (a) and (c) go \emph{down}. Suppose you had only the first two readings of each. Could
        you tell the two families apart? Say what the smallest amount of extra information would be
        that settles it, and what you would check.

\ansline{No --- any two points fit both a linear model and an exponential one, so two readings can never decide.}\\
\ansline{One more reading, at the same equal step, settles it: subtract to get the second difference and divide}\\
\ansline{to get the second ratio. Constant difference means linear; constant ratio means exponential.}\\
\end{enumerate}
\end{practicebox}

\end{document}
